% --- BẮT ĐẦU BẢNG ---

    

% Sử dụng cột c (canh giữa) cho STT và cột p (đoạn văn có kích thước cố định) cho nội dung để chữ tự xuống dòng
\begin{longtable}{|c|p{0.85\textwidth}|}
        \caption{Danh sách yêu cầu chức năng của hệ thống GeoAI} \label{tab:yeu_cau_chuc_nang} \\ % Thêm dòng này

\hline
\rowcolor{gray!25} % Tô màu xám nhạt cho hàng tiêu đề
\textbf{STT} & \centering\arraybackslash\textbf{Yêu cầu chức năng} \\ \hline
\endfirsthead

% Định nghĩa phần đầu của các trang tiếp theo (chỉ có đường kẻ ngang, KHÔNG lặp lại tiêu đề)
\hline
\endhead

% Định nghĩa phần cuối của các trang (khi bị ngắt trang)
\hline
\endfoot

% Định nghĩa phần cuối cùng của toàn bộ bảng
\hline
\endlastfoot

% --- NỘI DUNG BẢNG ---

\textbf{I} & \textbf{Chức năng hệ thống} \\ \hline
1 & Quản lý danh mục Cấu hình hệ thống \\ \hline
2 & Quản lý danh mục \gls{api} Key \\ \hline
3 & Quản lý danh mục Nhóm quyền \\ \hline
4 & Quản lý danh mục Chức năng \\ \hline
5 & Log \gls{api} (các phần mềm khác sử dụng \gls{api} của hệ thống) \\ \hline
6 & Nhật ký sử dụng hệ thống \\ \hline
7 & Sao lưu dữ liệu \\ \hline
8 & Phân quyền cho người dùng \\ \hline

\textbf{II} & \textbf{Chức năng \gls{webgis} -- Hiển thị bản đồ} \\ \hline
9 & Hiển thị bản đồ đô thị nền (OSM / Google Maps / VN Map), zoom 1--22 \\ \hline
10 & Quản lý lớp dữ liệu: bật/tắt lớp, kéo thả thứ tự, điều chỉnh độ trong suốt \\ \hline
11 & Hiển thị tài sản trên bản đồ với icon phân loại và popup thông tin khi click \\ \hline
12 & Tìm kiếm địa chỉ, vị trí và tài sản theo mã/tên; tự động zoom đến kết quả \\ \hline
13 & Lọc hiển thị tài sản theo loại, trạng thái, quận/phường, ngày cập nhật \\ \hline
14 & Vẽ và chỉnh sửa đối tượng không gian (điểm, đường, vùng) trực tiếp trên bản đồ \\ \hline
15 & Đo khoảng cách và diện tích trực tiếp trên bản đồ \\ \hline
16 & Xuất bản đồ sang PNG/PDF kèm legend, tỷ lệ; chia sẻ bản đồ qua link URL \\ \hline

\textbf{III} & \textbf{Chức năng Quản lý Tài sản} \\ \hline
17 & Thêm, sửa, xóa tài sản với đầy đủ thông tin thuộc tính, tọa độ và hình ảnh \\ \hline
18 & Quản lý hồ sơ tài sản: lịch sử bảo trì, tài liệu kỹ thuật, trạng thái và giá trị \\ \hline
19 & Quản lý lịch bảo trì định kỳ, cảnh báo đến hạn và ghi nhận kết quả thực hiện \\ \hline
20 & Import tài sản từ Excel/CSV/Shapefile; export sang Excel/GeoJSON/Shapefile \\ \hline
21 & Dashboard tổng quan tình trạng tài sản theo loại, trạng thái, đơn vị hành chính \\ \hline

\textbf{IV} & \textbf{Chức năng \gls{ai} -- Truy vấn thông minh} \\ \hline
22 & Truy vấn tài sản bằng ngôn ngữ tự nhiên tiếng Việt, trả kết quả trên bản đồ và bảng dữ liệu \\ \hline
23 & Hiển thị và cho phép chỉnh sửa câu SQL được sinh ra từ câu hỏi tự nhiên \\ \hline
24 & \gls{ai} đề xuất kế hoạch bảo trì dự đoán dựa trên lịch sử và tuổi thọ tài sản \\ \hline
25 & Nhận dạng tình trạng tài sản từ ảnh chụp hiện trường (hỏng hóc, vết nứt, gỉ sét) \\ \hline
26 & Tự động sinh báo cáo tình trạng tài sản bằng ngôn ngữ tự nhiên tiếng Việt, xuất PDF \\ \hline

\textbf{V} & \textbf{Chức năng Phân tích Không gian} \\ \hline
27 & Phân tích mật độ và phân bố tài sản theo không gian (heatmap) \\ \hline
28 & Phân tích vùng đệm (buffer) xung quanh tài sản, thống kê tài sản trong vùng \\ \hline
29 & Tối ưu tuyến đường bảo trì, tính lộ trình ngắn nhất cho danh sách tài sản cần xử lý \\ \hline
30 & Thống kê tài sản theo đơn vị hành chính: phường, quận, thành phố; bản đồ choropleth \\ 
\end{longtable}
